% One style for presentation fragments, worksheets, Markdown and external files.
% Native listings remain page-breakable: no float, minipage or enclosing box.
\definecolor{codebackground}{gray}{0.965}
\definecolor{codekeyword}{RGB}{35,61,83}
\definecolor{codestring}{RGB}{103,75,55}
\definecolor{codecomment}{gray}{0.40}
\definecolor{codenumber}{gray}{0.48}
\lstdefinelanguage{bookjson}{
  morestring=[b]",
  morekeywords={true,false,null},
  sensitive=true
}
% listings' default character table ends at U+00FF. Under XeTeX, characters
% above that range otherwise bypass its token buffer and can move within words
% (e.g. "Poslať" becomes "ťPosla"). Extend the table using its own processors,
% locally for each listing; retain the literal glyphs, including in strings.
\makeatletter
\begingroup
\catcode0=12\relax
\g@addto@macro\lst@DefEC{%
  \lst@CCECUse\lst@ProcessLetter čďěĺľňŕřšťůžČĎĚĹĽŇŔŘŠŤŮŽ^^00%
  \lst@CCECUse\lst@ProcessOther —€→←─│┌┐└┘├┤┬┴┼^^00%
}
\endgroup
\makeatother
\lstdefinestyle{bookcode}{
  language={},
  basicstyle=\small\ttfamily\color{black},
  keywordstyle=\bfseries\color{codekeyword},
  commentstyle=\itshape\color{codecomment},
  stringstyle=\color{codestring},
  columns=fullflexible,
  keepspaces=true,
  tabsize=4,
  showspaces=false,
  showtabs=false,
  showstringspaces=false,
  upquote=true,
  extendedchars=true,
  numbers=left,
  numberstyle=\scriptsize\ttfamily\color{codenumber},
  numbersep=9pt,
  firstnumber=1,
  stepnumber=1,
  numberblanklines=true,
  breaklines=true,
  breakatwhitespace=false,
  breakautoindent=true,
  breakindent=1em,
  backgroundcolor=\color{codebackground},
  frame=single,
  framerule=0pt,
  frameround=ffff,
  framesep=6pt,
  % The frame starts at the surrounding left margin. Its background includes
  % the number gutter, with 6pt padding outside it and beside the last column.
  xleftmargin=30pt,
  xrightmargin=6pt,
  framexleftmargin=24pt,
  framexrightmargin=0pt,
  aboveskip=0.8\baselineskip,
  belowskip=0.8\baselineskip,
  captionpos=t
}
\lstset{style=bookcode}
\lstnewenvironment{Code}[1][]{\lstset{style=bookcode,#1}}{}
